\documentclass[10pt]{article}

\usepackage[a4paper,margin=0.9in]{geometry}
\usepackage{microtype}
\usepackage{lmodern}
\usepackage[T1]{fontenc}
\usepackage[utf8]{inputenc}
\usepackage{titlesec}
\usepackage{enumitem}
\setlist{nosep, leftmargin=*}
\titlespacing*{\section}{0pt}{1.4ex}{0.6ex}
\titlespacing*{\subsection}{0pt}{1.0ex}{0.4ex}
\titleformat{\section}{\large\bfseries}{\thesection}{0.5em}{}
\titleformat{\subsection}{\normalsize\bfseries}{\thesubsection}{0.5em}{}

\usepackage{amsmath,amssymb,bm,mathtools}
\DeclarePairedDelimiter{\norm}{\lVert}{\rVert}
\DeclareMathOperator*{\argmin}{arg\,min}
\DeclareMathOperator{\STFT}{STFT}
\newcommand{\R}{\mathbb{R}}
\newcommand{\Mb}{\bm{M}}
\newcommand{\xb}{\bm{x}}
\newcommand{\yb}{\bm{y}}
\newcommand{\D}{\bm{D}}
\newcommand{\alphab}{\bm{\alpha}}

\usepackage{graphicx}
\usepackage{booktabs}
\usepackage{longtable}
\usepackage{caption}
\captionsetup{font=small, labelfont=bf, skip=4pt}
\usepackage[hidelinks]{hyperref}
\usepackage{xcolor}
\definecolor{mygreen}{HTML}{0B4F2F}
\definecolor{codebg}{HTML}{F6F6F6}
\definecolor{codeborder}{HTML}{D0D0D0}
\usepackage{tcolorbox}
\tcbset{colback=mygreen!4, colframe=mygreen!50, boxrule=0.4pt, arc=2pt,
        left=4pt, right=4pt, top=2pt, bottom=2pt}

\usepackage{listings}
\lstset{
  basicstyle=\ttfamily\footnotesize,
  backgroundcolor=\color{codebg},
  frame=single,
  rulecolor=\color{codeborder},
  framesep=4pt,
  framerule=0.4pt,
  numbers=left,
  numberstyle=\tiny\color{gray},
  numbersep=6pt,
  keywordstyle=\color{mygreen}\bfseries,
  stringstyle=\color{orange!80!black},
  commentstyle=\color{gray}\itshape,
  showstringspaces=false,
  breaklines=true,
  language=Python,
}

\title{\vspace{-2.0em}\textbf{Implementation Plan}\\[0.3em]
\large Sparse, Low-Rank, and Deep: Music Source Separation\\
\normalsize ENGS 109 (Spring 2026) Companion Document}
\author{Taka Khoo}
\date{\vspace{-1em}\today}

\begin{document}
\maketitle
\vspace{-2em}

\begin{tcolorbox}
\textbf{Purpose.}\;
The operational sibling of the proposal.  It answers \emph{how} the
project gets built, day by day, with playable audio at every
checkpoint.  Stage A already runs end-to-end (see
\texttt{audio\_demos/out/day0\_*.wav}); the rest of this document is
the path from there to the final report.
\end{tcolorbox}

% =====================================================================
\section{Sprint at a glance}
Each row below produces \emph{playable audio} that gets saved into
\texttt{audio\_demos/out/} and gets a row in the final report's
audio-examples table.  No deliverable is purely abstract.

\medskip
\noindent\small
\begin{tabular}{@{}p{1.2cm} p{4.4cm} p{4.0cm} p{4.6cm}@{}}
\toprule
\textbf{Day} & \textbf{What I build} & \textbf{Course tie} & \textbf{Audio output} \\
\midrule
0 (done) & Inexact-ALM Robust PCA solver, applied to MODULO mix &
L13 RPCA $(P_{R_1})$ &
\texttt{day0\_\{mix,harmonic,percussive\}.wav} \\[2pt]
1 & MUSDB18-HQ loader; render mix + 4 stems for first 5 train songs &
L1 motivation, L3 sampling &
\texttt{day1\_musdb\_*.wav} (15 files) \\[2pt]
2 & Stage A applied to MUSDB; museval SDR baseline number &
L13 + evaluation harness &
\texttt{day2\_stageA\_song\{1..5\}.wav} \\[2pt]
3 & Train $\D_v, \D_i$ NMF dictionaries on 20-song subset &
L15 Lee-Seung multiplicative updates &
\texttt{day3\_nmf\_atoms.pdf} (visual atom catalog) \\[2pt]
4 & Stage B inference on 5 test songs; per-source WAVs &
L15 separation via concatenated dictionary &
\texttt{day4\_stageB\_*.wav} (10 files) \\[2pt]
5 & K-SVD dictionaries per source, OMP-based code &
L14 K-SVD, PS3 OMP reused &
\texttt{day5\_ksvd\_atoms.pdf} \\[2pt]
6 & Stage C SRC inference; produce binary masks &
L8 SRC residual rule &
\texttt{day6\_stageC\_*.wav} \\[2pt]
7-10 & SparseNet skeleton in PyTorch, overfit one song &
L16 composite module + L10 FISTA backward &
\texttt{day10\_overfit\_one\_song.wav} \\[2pt]
11-18 & Full SparseNet training on MUSDB18-HQ train split &
L16 end-to-end backprop &
\texttt{day18\_stageD\_test\_set.zip} \\[2pt]
19-22 & MMV stereo refinement (SOMP); add joint-sparsity post-pass &
L12 simultaneous OMP &
\texttt{day22\_stageE\_stereo\_*.wav} \\[2pt]
23-28 & Full evaluation harness vs Demucs, Spleeter, ElevenLabs;
ablations on $K$, $S$, $\lambda$ &
L6/L7 theoretical bounds applied to results &
\texttt{day28\_eval\_report.pdf} \\[2pt]
29-35 & MUSHRA-style listening test with 8 raters &
listener-grounded evaluation &
\texttt{day35\_listening\_study.csv} \\[2pt]
36-42 & Final report write-up + code release &
$-$ &
\texttt{final\_report.pdf} \\
\bottomrule
\end{tabular}

\paragraph{Sprint rule.} If a day slips, drop the audio output to a
shorter clip but never drop the audio entirely.  The day-by-day audio
trail \emph{is} the project log.

\section{Executive summary}
Total fresh code budget: roughly $1500$ lines of Python.  Of the five
pipeline stages, four reuse code I wrote during PS1--PS8; only
Stage~D (SparseNet) is a genuinely new implementation in PyTorch.
No new mathematics is required beyond what ENGS~109 covered.  Hardware:
MacBook Pro M3 for Stages A--C and E; Dartmouth Discovery cluster
(one A100) for Stage~D.  Day-0 RPCA already runs in $9.2$~s on the
MacBook with no GPU.

% =====================================================================
\section{Repository layout}
\begin{lstlisting}[language=bash]
engs109-music-separation/
|-- README.md
|-- environment.yml
|-- pyproject.toml
|-- data/
|   |-- musdb18hq/             # Symlink to /scratch dataset.
|   `-- cache/                 # Pre-computed STFTs (.npz).
|-- src/separation/
|   |-- __init__.py
|   |-- io.py                  # WAV loading, MUSDB18 iterator.
|   |-- features.py            # STFT, mel, CQT, inverse.
|   |-- masks.py               # Mask <-> waveform helpers.
|   |-- metrics.py             # museval + mir_eval wrappers.
|   |-- stage_a_rpca.py        # Inexact-ALM RPCA (from PS8).
|   |-- stage_b_nmf.py         # NMF + class-conditional dictionaries.
|   |-- stage_c_ksvd.py        # K-SVD + SRC (from PS3, PS7).
|   |-- stage_d_scn.py         # SparseNet (new, PyTorch).
|   |-- stage_e_mmv.py         # SOMP for stereo (extends PS3 OMP).
|   `-- train.py               # Stage D training entry point.
|-- notebooks/
|   |-- 00_data_smoke_test.ipynb
|   |-- 01_stage_a_rpca.ipynb
|   |-- 02_stage_b_nmf.ipynb
|   |-- 03_stage_c_ksvd.ipynb
|   |-- 04_stage_d_scn_overfit.ipynb
|   |-- 05_evaluation.ipynb
|   `-- 06_listening_study.ipynb
|-- experiments/               # Configs in YAML.
|-- reports/                   # Final figures, .tex.
|-- tests/                     # Unit tests for solvers.
`-- audio_examples/            # Curated test clips for the report.
\end{lstlisting}

% =====================================================================
\section{Environment setup}
\subsection*{environment.yml}
\begin{lstlisting}[language={}]
name: engs109-mss
channels: [conda-forge, pytorch]
dependencies:
  - python=3.11
  - numpy=2.0
  - scipy=1.13
  - matplotlib=3.9
  - scikit-learn=1.5
  - librosa=0.10
  - soundfile=0.12
  - ffmpeg=6.1
  - pytorch=2.3
  - cvxpy=1.5            # used in PS7, kept for BP/LP baselines
  - jupyterlab=4.2
  - pip
  - pip:
      - musdb==0.4.0
      - museval==0.4.1
      - mir_eval==0.7
      - openunmix==1.2.1
      - demucs==4.0.1
\end{lstlisting}
The same file is checked into the repository so that \texttt{conda env
create -f environment.yml} reconstructs everything bit-identically.

\subsection*{Hardware}
Local dev: MacBook Pro M3, 36~GB unified memory; runs Stages A--C and
all evaluation comfortably.  Stage~D training: Dartmouth Discovery
cluster, one NVIDIA A100 (40~GB) for an estimated $40$--$60$ wall-clock
hours over the 10-week window.

% =====================================================================
\section{Data pipeline}
\subsection*{Dataset acquisition}
MUSDB18-HQ ships as one \texttt{.tar.gz} ($\approx\!22$~GB) at
\href{https://zenodo.org/record/3338373}{zenodo.org/record/3338373}.
Extract to \texttt{/scratch/musdb18hq/} (Discovery) or
\texttt{\$HOME/data/musdb18hq/} (local) and symlink to
\texttt{data/musdb18hq/}.  Standard splits: $100$ training, $50$ test
songs.  I will reserve $14$ songs from train as validation
(\texttt{musdb.DB(subsets="train", split="valid")}).

\subsection*{Feature extraction}
Reuse the STFT helper I already wrote for \texttt{MozartServe/audio2midi/cqt.py}:
\begin{lstlisting}
# src/separation/features.py
import librosa
import numpy as np

SR, N_FFT, HOP, WIN = 22050, 2048, 512, 2048

def stft(y: np.ndarray) -> np.ndarray:
    """Complex STFT; rows = freq bins, cols = time frames."""
    return librosa.stft(y, n_fft=N_FFT, hop_length=HOP, win_length=WIN)

def istft(Y: np.ndarray) -> np.ndarray:
    return librosa.istft(Y, hop_length=HOP, win_length=WIN)

def magphase(Y):       # complex -> (magnitude, phase)
    return np.abs(Y), np.angle(Y)

def apply_mask(M: np.ndarray, Y: np.ndarray) -> np.ndarray:
    """Soft mask M in [0,1] on magnitude, keep mixture phase."""
    return istft(M * np.abs(Y) * np.exp(1j * np.angle(Y)))
\end{lstlisting}
For the SparseNet stage I additionally use mel-band spectrograms
($64$ mels) so input patches are small enough for FISTA inference to stay
under $50$~ms per frame.

\subsection*{Augmentation (training-only)}
On-the-fly stem remixing: pick four random stems from different songs
(\texttt{vocals}, \texttt{drums}, \texttt{bass}, \texttt{other}), apply
a random gain $g\sim\mathcal{U}(-6, 6)$~dB to each, and sum.  This is
the standard SiSEC-2018 augmentation and roughly $20\times$'s the
effective dataset size.  Implemented in a custom
\texttt{torch.utils.data.IterableDataset}.

% =====================================================================
\section{Code reuse map}
Table~\ref{tab:reuse} keys every component back to existing problem-set
code.  I will refactor each PS notebook into a single library module so
nothing gets re-derived from scratch.

\begin{longtable}{@{}p{4.0cm} p{4.6cm} p{6.8cm}@{}}
\caption{Existing problem-set code that becomes the spine of each stage.}
\label{tab:reuse}\\
\toprule
\textbf{Module} & \textbf{Where it comes from} & \textbf{What needs changing} \\
\midrule
\endfirsthead
\toprule
\textbf{Module} & \textbf{Where it comes from} & \textbf{What needs changing} \\
\midrule
\endhead
OMP solver  &
PS3 \texttt{omp(A, y, S)} &
Vectorise over time frames; add MMV variant for Stage~E. \\
ISTA/FISTA  &
PS5 \texttt{fista(A, y, lam)} &
Wrap as a \texttt{torch.autograd.Function} for Stage~D unrolling. \\
Basis Pursuit / LASSO  &
PS7 \texttt{cvxpy} formulations &
Only used as a sanity check vs FISTA. \\
SRC classifier  &
PS7 face-recognition notebook &
Replace class labels with $\{$vocal, instrumental$\}$. \\
K-SVD &
PS3 supplementary &
Add a per-source training loop. \\
Inexact-ALM nuclear norm &
PS8 matrix-completion notebook &
Add the sparse $\bm{E}$ update for full RPCA. \\
NMF multiplicative updates &
\emph{New} (sklearn baseline first) &
Implement KL-divergence variant from L15 slides. \\
Stage D SparseNet &
\emph{Fully new}, PyTorch &
4 composite modules; see \S\ref{sec:scn-code}. \\
Evaluation harness &
\emph{New}, thin wrapper &
Calls \texttt{museval.eval\_mus\_track}. \\
\bottomrule
\end{longtable}

% =====================================================================
\section{Per-stage implementation}

\subsection{Stage A. Robust PCA (4 days)}
The full algorithm is the inexact ALM of Lin--Chen--Ma (2010), already
implemented in the figure-generation script for the proposal and verified
to recover the synthetic $60\!\times\!80$ rank-3 matrix to $10^{-7}$ in
$32$ iterations.

\begin{lstlisting}
# src/separation/stage_a_rpca.py
import numpy as np

def shrink(X, tau):
    return np.sign(X) * np.maximum(np.abs(X) - tau, 0.0)

def svt(X, tau):
    U, s, Vt = np.linalg.svd(X, full_matrices=False)
    s = np.maximum(s - tau, 0.0)
    return U @ np.diag(s) @ Vt

def rpca(Y, lam=None, tol=1e-7, max_iter=200):
    n1, n2 = Y.shape
    lam = 1.0 / np.sqrt(max(n1, n2)) if lam is None else lam
    norm_two = np.linalg.norm(Y, 2)
    Y2 = Y / max(norm_two, np.linalg.norm(Y, np.inf) / lam)
    mu, rho = 1.25 / norm_two, 1.5
    L, E = np.zeros_like(Y), np.zeros_like(Y)
    for k in range(max_iter):
        L = svt(Y - E + Y2 / mu, 1.0 / mu)
        E = shrink(Y - L + Y2 / mu, lam / mu)
        Y2 = Y2 + mu * (Y - L - E)
        mu = mu * rho
        if np.linalg.norm(Y - L - E, "fro") / np.linalg.norm(Y, "fro") < tol:
            break
    return L, E
\end{lstlisting}
Applied to magnitude spectrograms $\Mb_Y\!\in\!\R_{\ge 0}^{F\times T}$,
$L$ is the harmonic stem estimate, $E$ the percussive one.  Wiener
masking gives soft masks:
$\Mb_h = L^2 / (L^2 + E^2),\;\Mb_p = E^2 / (L^2 + E^2)$.

\subsection{Stage B. NMF with class-conditional dictionaries (5 days)}
Training: separate Lee--Seung loops on vocal-only and instrumental-only
training stems, $K_v\!=\!K_i\!=\!64$, $400$ iters, KL divergence.
sklearn's \texttt{NMF(beta\_loss="kullback-leibler", solver="mu")} works
out of the box.  Test-time decomposition holds $\bm{D}$ fixed and solves
only for $\bm{H}$.

\begin{lstlisting}
# src/separation/stage_b_nmf.py
from sklearn.decomposition import NMF
import numpy as np

def train_dictionary(M_train, K=64, max_iter=400):
    """M_train: (F, T_all) nonnegative magnitude spectrogram."""
    nmf = NMF(n_components=K, init="nndsvda", solver="mu",
              beta_loss="kullback-leibler", max_iter=max_iter)
    H = nmf.fit_transform(M_train.T).T   # H: (K, T_all)
    D = nmf.components_.T                 # D: (F, K)
    return D, H

def separate(M_mix, D_v, D_i, beta=1e-2, max_iter=200):
    """Joint factorisation with concatenated dictionary."""
    F, T = M_mix.shape
    D = np.hstack([D_v, D_i])
    Kv, Ki = D_v.shape[1], D_i.shape[1]
    H = np.random.rand(Kv + Ki, T) + 1e-3
    for _ in range(max_iter):
        num = D.T @ M_mix
        den = D.T @ D @ H + beta + 1e-9
        H *= num / den
    Mv_hat = D_v @ H[:Kv]
    Mi_hat = D_i @ H[Kv:]
    mask_v = (Mv_hat ** 2) / (Mv_hat ** 2 + Mi_hat ** 2 + 1e-9)
    return mask_v
\end{lstlisting}
Output \texttt{mask\_v} multiplies the complex STFT then inverts.
Two ablations to run: (a) Frobenius vs KL loss, (b) $K\in\{32,64,128\}$.

\subsection{Stage C. K-SVD + SRC (5 days)}
Train one K-SVD dictionary per source from $4{,}096$ randomly-sampled
spectrogram frames each.  Per-frame sparsity $S\!=\!10$, dictionary size
$K\!=\!256$, $80$ outer iterations.  Implementation reuses my PS3 OMP
and adds the column-by-column SVD update from L14:

\begin{lstlisting}
# src/separation/stage_c_ksvd.py
import numpy as np

def omp(D, y, S):
    """Reused from PS3; returns S-sparse x with y ~= D x."""
    residual = y.copy()
    support, x = [], np.zeros(D.shape[1])
    for _ in range(S):
        idx = int(np.argmax(np.abs(D.T @ residual)))
        if idx in support:
            break
        support.append(idx)
        x_s, *_ = np.linalg.lstsq(D[:, support], y, rcond=None)
        x[:] = 0
        x[support] = x_s
        residual = y - D @ x
    return x

def ksvd(Y, K, S, n_iter=80):
    F, L = Y.shape
    D = Y[:, np.random.choice(L, K, replace=False)]
    D /= np.linalg.norm(D, axis=0, keepdims=True) + 1e-9
    X = np.zeros((K, L))
    for it in range(n_iter):
        for j in range(L):
            X[:, j] = omp(D, Y[:, j], S)
        for k in range(K):
            wk = np.flatnonzero(X[k])
            if len(wk) == 0:
                continue
            Ek = Y[:, wk] - D @ X[:, wk] + np.outer(D[:, k], X[k, wk])
            U, s, Vt = np.linalg.svd(Ek, full_matrices=False)
            D[:, k] = U[:, 0]
            X[k, wk] = s[0] * Vt[0]
    return D, X
\end{lstlisting}

At test time, build the union dictionary $\bm{D}\!=\![\bm{D}_v\,|\,\bm{D}_i]$,
run OMP per frame, then assign each non-zero atom to its source.
Per-bin source assignment follows the L8 residual rule.

\subsection{Stage D. SparseNet mask predictor (3 weeks)}
\label{sec:scn-code}
The genuinely new code.  The composite sparse coding module of L16 is
implemented as a PyTorch \texttt{Module} whose forward pass unrolls
$T_\text{f}\!=\!30$ FISTA iterations for both the upsampling and
downsampling sparse codes; gradients flow back via PyTorch autograd.

\begin{lstlisting}
# src/separation/stage_d_scn.py  (sketch)
import torch
import torch.nn as nn

def fista_nn(D, y, lam1, lam2, T=30):
    """Non-negative FISTA for (1/2)||y - D a||^2 + lam1||a||_1 + (lam2/2)||a||_2^2."""
    L = torch.linalg.matrix_norm(D, ord=2) ** 2 + lam2
    a = torch.zeros(D.shape[1], device=D.device)
    z, t = a.clone(), 1.0
    for _ in range(T):
        grad = D.T @ (D @ z - y) + lam2 * z
        a_new = torch.clamp_min(z - grad / L - lam1 / L, 0.0)
        t_new = (1 + (1 + 4 * t ** 2) ** 0.5) / 2
        z = a_new + ((t - 1) / t_new) * (a_new - a)
        a, t = a_new, t_new
    return a

class CompositeModule(nn.Module):
    def __init__(self, n_in, n_up, n_down, lam1=0.05, lam2=0.05):
        super().__init__()
        self.D1 = nn.Parameter(torch.randn(n_in, n_up) * 0.05)
        self.D2 = nn.Parameter(torch.randn(n_up, n_down) * 0.05)
        self.lam1 = nn.Parameter(torch.tensor(lam1))
        self.lam2 = nn.Parameter(torch.tensor(lam2))
    def forward(self, x):                            # x: (B, n_in)
        a1 = torch.stack([fista_nn(self.D1, xi, self.lam1, self.lam2)
                          for xi in x])
        a2 = torch.stack([fista_nn(self.D2, ai, self.lam1, self.lam2)
                          for ai in a1])
        return a2                                    # (B, n_down)

class SparseNet(nn.Module):
    def __init__(self, F=64, L=4, n_up=128, n_down=64):
        super().__init__()
        dims = [F] + [n_down] * L
        self.modules_ = nn.ModuleList(
            CompositeModule(dims[i], n_up, dims[i + 1]) for i in range(L))
        self.head = nn.Conv2d(n_down, 1, kernel_size=1)
    def forward(self, X):                            # X: (B, F, T)
        h = X.transpose(1, 2).reshape(-1, X.shape[1])   # (B*T, F)
        for mod in self.modules_:
            h = mod(h)
        h = h.reshape(X.shape[0], X.shape[2], -1).transpose(1, 2)
        h = h.unsqueeze(1)                              # (B, 1, n_down, T)
        return torch.sigmoid(self.head(h)).squeeze(1)  # (B, F, T) mask
\end{lstlisting}

Two implementation details that are not optional:
\begin{enumerate}
\item \textbf{Batch the FISTA.} The naive per-sample loop in
\texttt{forward} above is for readability; the production code batches
into a single $(B,\,N_u)$ matmul per iteration so a 30-step inference is
one fused kernel call.
\item \textbf{Pre-compute Lipschitz constants.} L16 explicitly notes
``Matrix inversion is pre-computed.''  At each gradient step I cache
$\|\bm{D}^{(\ell)}\|_2^2$ via power iteration once per epoch rather than
recomputing each forward pass.
\end{enumerate}

\paragraph{Training loop.}
\begin{lstlisting}
# src/separation/train.py
import torch
from separation.stage_d_scn import SparseNet
from separation.features import stft, magphase

net = SparseNet(F=64, L=4).cuda()
opt = torch.optim.AdamW(net.parameters(), lr=3e-4, weight_decay=1e-5)

for epoch in range(100):
    for mix_wav, voc_wav in loader:
        Ymix, Yvoc = stft(mix_wav), stft(voc_wav)
        Mmix, Pmix = magphase(Ymix)
        Mvoc, _    = magphase(Yvoc)
        mask = net(Mmix.cuda())
        recon = mask * torch.as_tensor(Mmix).cuda()
        loss = (recon - torch.as_tensor(Mvoc).cuda()).abs().mean()
        loss = loss + 1e-3 * sum(p.abs().mean() for p in net.parameters()
                                 if p.requires_grad)
        opt.zero_grad(); loss.backward(); opt.step()
\end{lstlisting}

\paragraph{Hyperparameters (fixed at the start, then frozen).}
$L_\text{f}{=}30$ FISTA steps, $\lambda_1\!=\!\lambda_2\!=\!0.05$ initial,
4 composite modules, fat dictionary width $N_u\!=\!128$, tall dictionary
width $N_d\!=\!64$, batch $B\!=\!16$ song-snippets of 4~s, learning rate
$3\!\times\!10^{-4}$ AdamW with cosine decay, 100 epochs.  Mixed
precision (\texttt{torch.amp}) on the A100.

\subsection{Stage E. SOMP for stereo (3 days)}
The simultaneous OMP of L12 needs only a one-line change to the PS3 OMP:
the scoring step becomes a row-norm of the residual matrix.

\begin{lstlisting}
def somp(D, Y, S, q=2):
    """Y: (F, K_meas).  Returns row-sparse X: (atoms, K_meas)."""
    R = Y.copy()
    support = []
    for _ in range(S):
        scores = np.linalg.norm(D.T @ R, ord=q, axis=1)
        idx = int(np.argmax(scores))
        if idx in support:
            break
        support.append(idx)
        Xs, *_ = np.linalg.lstsq(D[:, support], Y, rcond=None)
        R = Y - D[:, support] @ Xs
    X = np.zeros((D.shape[1], Y.shape[1]))
    X[support] = Xs
    return X
\end{lstlisting}
At inference, stack $\yb_L,\yb_R$ as the $2$-column $\bm{Y}$ and call
\texttt{somp(D, Y, S=10)}.  The shared row support is the joint sparse
code; per-channel reconstructions then go through the standard inverse
STFT with their own phases.

% =====================================================================
\section{Training schedule}
\begin{tabular}{@{}lll@{}}
\toprule
\textbf{Stage} & \textbf{Training data} & \textbf{Wall-clock estimate} \\
\midrule
A. RPCA           & none (no training)            & instantaneous \\
B. NMF            & $1000$ vocal + $1000$ instr.\ frames & 8 min CPU \\
C. K-SVD          & $4096$ frames per source       & 25 min CPU \\
D. SparseNet      & MUSDB18 train, augmented       & 45--60 h A100 \\
E. SOMP           & none (uses Stage C dict.)      & instantaneous \\
\bottomrule
\end{tabular}

% =====================================================================
\section{Evaluation harness}
\begin{lstlisting}
# src/separation/metrics.py
import museval, musdb
def evaluate_on_musdb(separator_fn, output_dir):
    db = musdb.DB(root="data/musdb18hq", subsets=["test"])
    results = museval.EvalStore()
    for track in db:
        estimates = separator_fn(track.audio, track.rate)
        scores = museval.eval_mus_track(track, estimates,
                                        output_dir=output_dir)
        results.add_track(scores)
    return results
\end{lstlisting}
\texttt{separator\_fn} returns a dict
\texttt{\{"vocals": ndarray, "accompaniment": ndarray\}} at the original
sample rate.  Reports median SDR/SI-SDR/SIR/SAR per stem with
$25$th--$75$th percentile bars, matching the SiSEC convention.

% =====================================================================
\section{Notebook plan (mirroring PS format)}
Each notebook follows the PS template I used all term: a header cell with
the problem statement, math derivation in markdown, a single solver
function, a synthetic verification block, then the real-data block.

\begin{itemize}
\item \textbf{00\_data\_smoke\_test.ipynb}: confirm MUSDB18-HQ loads,
plot the first track's STFT/CQT/mel of each stem.
\item \textbf{01\_stage\_a\_rpca.ipynb}: synthetic verification (as in
the proposal Fig.~2), then apply to a MUSDB track; render
$X$/$E$ spectrograms; tabulate SDR.
\item \textbf{02\_stage\_b\_nmf.ipynb}: train $\bm{D}_v$ and $\bm{D}_i$
on training stems; show learned atoms; report SDR with KL vs Frobenius.
\item \textbf{03\_stage\_c\_ksvd.ipynb}: K-SVD training on synthetic
sparse signals first (sanity), then on spectrogram frames; SRC residual
visualisation.
\item \textbf{04\_stage\_d\_scn\_overfit.ipynb}: deliberately overfit
SparseNet to a single MUSDB track to verify backward pass; check that
intermediate codes are actually sparse ($>\!95\%$ zeros).
\item \textbf{05\_evaluation.ipynb}: run all five stages on the
$50$-song test set; produce the headline SDR/SI-SDR boxplots and the
per-track CSVs.
\item \textbf{06\_listening\_study.ipynb}: pick 20 clips, anonymise,
collect MUSHRA scores from $8$ raters; box-plot the results.
\end{itemize}

% =====================================================================
\section{Risk register and fallbacks}

\begin{longtable}{@{}p{4.0cm} p{4.5cm} p{6.5cm}@{}}
\toprule
\textbf{Risk} & \textbf{Detection} & \textbf{Fallback} \\
\midrule
\endhead
SparseNet too slow per iter & $>\!200$~ms/frame on A100 &
Drop to $L\!=\!2$ composite modules, replace mel input with the first $32$
NMF activations from Stage~B (effectively shrinks $n_\text{in}$ from
$64$ to $32$). \\
SparseNet does not converge & validation SDR flat after $30$ epochs &
Warm-start $\bm{D}^{(1)}$ from K-SVD-trained dictionary (Stage~C);
linearly schedule $\lambda_1\colon 0.5 \!\to\! 0.05$. \\
MUSDB-HQ ground-truth alignment drifts &
Per-track SDR has bimodal distribution &
Run \texttt{museval} alignment compensation; document as known issue. \\
Listening-study response rate low &
$<\!4$ raters in week 9 &
Reduce sample count to $10$ clips; use myself + 2 Mozart~AI teammates
as fallback raters. \\
Discovery cluster queue blocks training &
no GPU available for $>\!48$ h &
Train Stage~D locally on M3 with reduced model: $L\!=\!2$, batch~$=4$,
$15$ epochs (about $20$ wall-clock hours). \\
\bottomrule
\end{longtable}

% =====================================================================
\section{Appendix A: Expected figure list}
For the final report (separate document from this plan):
\begin{itemize}
\item Fig.~1: pipeline diagram (already produced for proposal).
\item Fig.~2: spectrogram triptych mix/vocals/instrumental (already produced).
\item Fig.~3: chromagram + harmonic--percussive split for one
representative MUSDB song.
\item Fig.~4: RPCA decomposition of spectrogram with quantified
SDR per stage.
\item Fig.~5: NMF learned atoms $\bm{D}_v,\bm{D}_i$ as small spectral
heatmaps (analogue of the SCN feature-map figure on L16 slide~14).
\item Fig.~6: K-SVD atoms for both sources, with per-atom sparsity
histograms.
\item Fig.~7: SparseNet intermediate sparsity profile; demonstrating
the L16 ``$>\!95\%$ zeros'' regime.
\item Fig.~8: SDR boxplots on MUSDB18-HQ test set, all stages and
baselines.
\item Fig.~9: MUSHRA listening-study results.
\item Fig.~10: ablation table on $L$, $K$, $S$, $\lambda_1$.
\end{itemize}

% =====================================================================
\section{Appendix B: Connection to my Mozart~AI codebase}
\paragraph{What I will lift directly:}
\begin{itemize}
\item \texttt{MozartServe/audio2midi/cqt.py} provides constant-Q-transform
helpers that I will reuse for the CQT variant of Stage~B.
\item \texttt{MozartServe/audio2midi/midi.py} for ground-truth onset/offset
tagging when I want to qualify separation quality on monophonic
instrumental excerpts.
\item Mozart~AI's user-track manifest format (one JSON per project) gives
me a ready-made schema for storing the listening-study trials.
\end{itemize}

\paragraph{What goes back into Mozart~AI:}
If Stage~D outperforms ElevenLabs on vocal SDR (the success criterion),
I drop the trained checkpoint into a new
\texttt{MozartServe/separation/} package and replace the ElevenLabs
HTTP call in the ``Extract Stems'' button.  This closes the
$35$~dB SI-SDR gap from my MS thesis and removes one paid third-party
dependency from the product.

% =====================================================================
\section{Appendix C: What I am explicitly \emph{not} doing}
To keep this honest about scope:
\begin{itemize}
\item No multi-source separation beyond vocals vs accompaniment.  4-stem
(vocals/drums/bass/other) is a stretch goal only.
\item No time-domain models (Wave-U-Net, Conv-TasNet).  The whole point
of the project is to stay in the spectrogram CS regime.
\item No transformer architectures.  L16's SparseNet is the deep
architecture the course gave us; that is the one I am evaluating.
\item No real-time inference target.  Sample-level offline batch is
the only operating point.
\item No additional datasets beyond MUSDB18-HQ.  Cross-dataset
generalisation is left for future work.
\end{itemize}

\end{document}
